\notechapter{N-20260618}{First-Order and Higher-Order Differential Equations}

\section{Why the differential-equation part is one note}

The last part of the original long Markdown note covers first-order equations and higher-order linear equations.  These belong together because the central idea is the same: transform an equation for an unknown function into a form that can be integrated, linearized, or solved as an operator equation.

\section{Basic concepts}

A differential equation relates an unknown function to its derivatives.  The order is the highest derivative appearing.  A general first-order equation has the form
\[
  F(x,y,y')=0.
\]
A solution is a function $y(x)$ satisfying the equation on an interval.  A general solution contains arbitrary constants; an initial condition selects a particular solution.

\section{Separable equations}

A separable equation has
\[
  \frac{dy}{dx}=g(x)h(y).
\]
If $h(y)\ne0$, then
\[
  \frac{dy}{h(y)}=g(x)\,dx.
\]
Integrating gives an implicit solution.  One must also check equilibrium solutions coming from $h(y)=0$, because division by $h(y)$ can lose them.

\section{Homogeneous first-order equations}

A homogeneous first-order equation has the form
\[
  y'=F\left(\frac yx\right).
\]
Set $v=y/x$, so $y=vx$ and
\[
  y'=v+xv'.
\]
Then
\[
  v+x\frac{dv}{dx}=F(v),
\]
which is separable:
\[
  \frac{dv}{F(v)-v}=\frac{dx}{x}.
\]

Some equations become homogeneous after translating the origin.  If the linear expressions in $x,y$ meet at a point, shift that point to the origin.

\section{First-order linear equations}

A first-order linear equation is
\[
  y'+p(x)y=q(x).
\]
The integrating factor is
\[
  \mu(x)=e^{\int p(x)\,dx}.
\]
Multiplying the equation by $\mu$ gives
\[
  (\mu y)'=\mu q.
\]
Thus
\[
  y=e^{-\int p}
  \left(\int q e^{\int p}\,dx+C\right).
\]
The integrating factor is not a trick; it conjugates the operator $D+p(x)$ into the ordinary derivative $D$.

\section{Bernoulli equations}

A Bernoulli equation has
\[
  y'+p(x)y=q(x)y^n.
\]
For $n\ne0,1$, set
\[
  u=y^{1-n}.
\]
Then
\[
  u'=(1-n)y^{-n}y',
\]
and the equation becomes linear in $u$.  Again, check special solutions lost by division.

\section{Exact equations}

An equation
\[
  M(x,y)\,dx+N(x,y)\,dy=0
\]
is exact if there is a potential $\Phi$ with
\[
  d\Phi=M\,dx+N\,dy.
\]
The test is
\[
  M_y=N_x
\]
on a suitable simply connected region.  Then the solution is
\[
  \Phi(x,y)=C.
\]
If the equation is not exact, one may look for an integrating factor, often depending only on $x$ or only on $y$.

\section{Reducible higher-order equations}

If a second-order equation does not explicitly contain $y$, set $p=y'$ and reduce the order.  If it does not explicitly contain $x$, set $p=p(y)=y'$, so
\[
  y''=\frac{dp}{dx}=\frac{dp}{dy}\frac{dy}{dx}=p\frac{dp}{dy}.
\]
These substitutions are not random: they use a missing variable as a symmetry.

\section{Second-order linear equations}

A second-order linear equation has
\[
  y''+p(x)y'+q(x)y=f(x).
\]
The associated homogeneous equation has a solution space.  If $y_1,y_2$ are independent homogeneous solutions, then
\[
  y_h=C_1y_1+C_2y_2.
\]
A general nonhomogeneous solution is
\[
  y=y_h+y_p.
\]
This is the affine-space principle: one particular solution plus the kernel of a linear operator.

The Wronskian
\[
  W(y_1,y_2)=y_1y_2'-y_1'y_2
\]
measures independence.  If it is nonzero at one point for solutions of a linear homogeneous equation, the two solutions form a fundamental system.

\section{Constant-coefficient homogeneous equations}

For
\[
  y''+ay'+by=0,
\]
try $y=e^{rx}$.  The characteristic equation is
\[
  r^2+ar+b=0.
\]
Distinct real roots $r_1,r_2$ give
\[
  y=C_1e^{r_1x}+C_2e^{r_2x}.
\]
A repeated root $r$ gives
\[
  y=(C_1+C_2x)e^{rx}.
\]
Complex roots $\alpha\pm i\beta$ give
\[
  y=e^{\alpha x}(C_1\cos\beta x+C_2\sin\beta x).
\]

\section{Constant-coefficient nonhomogeneous equations}

For
\[
  L[y]=f(x),\qquad L=D^2+aD+b,
\]
the method of undetermined coefficients works when $f$ is built from polynomials, exponentials, sines, and cosines.  If the trial function overlaps with a homogeneous solution, multiply by a sufficient power of $x$.

Variation of parameters works more generally.  If $y_1,y_2$ are fundamental solutions, seek
\[
  y_p=u_1(x)y_1(x)+u_2(x)y_2(x).
\]
The auxiliary condition $u_1'y_1+u_2'y_2=0$ produces a solvable linear system for $u_1',u_2'$.

\section{Initial-value problems}

An initial-value problem specifies
\[
  y(x_0)=a,\qquad y'(x_0)=b.
\]
For regular linear equations, existence and uniqueness say that these data select exactly one solution.  This is conceptually important: a second-order equation needs two pieces of initial data because its solution space is two-dimensional.

\section{Functional-analytic viewpoint}

Let
\[
  L[y]=y''+p(x)y'+q(x)y.
\]
Then a differential equation is an operator equation
\[
  L[y]=f.
\]
The homogeneous solutions form $\ker L$.  A nonhomogeneous solution set, if nonempty, is an affine translate of $\ker L$.  Initial conditions are linear functionals that choose one element from that affine space.

This viewpoint explains why linear differential equations behave like linear algebra in function space.  The derivative operator replaces a matrix; the Wronskian replaces a determinant; fundamental solutions replace a basis; Green functions, when introduced later, are kernels for inverse operators.

\section{Differential equations as operators}

The elementary note classifies equations by solvable forms: separable, homogeneous, linear, Bernoulli, exact, constant-coefficient.  The higher-level unity is that a differential equation is an equation in a function space.  For example,
\[
  L[y]=f,\qquad L=\frac{d^2}{dx^2}+p(x)\frac d{dx}+q(x).
\]
The homogeneous solutions form the kernel $\ker L$.  A nonhomogeneous solution set, when nonempty, is an affine translate
\[
  y_p+\ker L.
\]
This is the same algebraic structure as a linear system $Ax=b$.

\section{Green functions}

For a linear operator $L$, a Green function $G(x,\xi)$ is a kernel that represents the inverse problem.  Formally,
\[
  L_xG(x,\xi)=\delta(x-\xi),
\]
where $\delta$ is the Dirac delta distribution.  Then a solution of
\[
  L[y]=f
\]
can often be written as
\[
  y(x)=\int G(x,\xi)f(\xi)\,d\xi
\]
with boundary conditions built into $G$.

This connects elementary particular-solution methods to integral operators.  Undetermined coefficients works for special forcing terms; Green functions describe the response to arbitrary forcing as a superposition of responses to point sources.

\section{Semigroups and evolution}

For first-order evolution equations, one often writes
\[
  \frac{du}{dt}=Au.
\]
If $A$ were a number, the solution would be $u(t)=e^{tA}u(0)$.  In a Banach or Hilbert space, $A$ may be an unbounded operator, and the solution is described by a semigroup
\[
  T(t)=e^{tA},\qquad T(t+s)=T(t)T(s).
\]
The heat equation is the model example: the Laplacian generates a smoothing semigroup.  The elementary exponential solution of constant-coefficient ODEs is the finite-dimensional shadow of this theory.

\section{Sturm--Liouville and Fourier expansion}

A second-order operator of the form
\[
  L[y]=-(p(x)y')'+q(x)y
\]
with suitable boundary conditions is often self-adjoint.  Its eigenvalue problem
\[
  L[y]=\lambda w(x)y
\]
produces orthogonal eigenfunctions.  Fourier sine and cosine series are special cases of this principle.

Thus the elementary method of solving $y''+\lambda y=0$ under boundary conditions is not an isolated trick.  It is the beginning of spectral theory for differential operators.
